\chapter{Volume constraints and CMC stability}\label{sec:cmc-stability}
Stationarity and stability depend on which variations are allowed. We distinguish critical points of area minus a prescribed multiple of volume from critical points of area on a fixed-volume class. The second-variation form is the same at a CMC hypersurface, but its domain of test functions is different. Sources are \cite{ZZCMC}, equations (1.2)--(1.3) and Definition 2.1, and the smooth variation calculations in \cite{HP99} and \cite{ChMin}, Section 23.

\begin{definition}[Infinitesimal volume-preserving variations]\label{def:volume-preserving-speed}
Let $\Omega$ be a smooth compact region in a Riemannian manifold, with nonempty smooth boundary $\Sigma$ and outward normal $\nu$. A normal speed $f\in C^\infty(\Sigma)$ preserves volume to first order when
\[
 \int_\Sigma f\,dA=0.
\]
An exactly volume-preserving variation instead satisfies $\operatorname{Vol}(\Omega_t)=\operatorname{Vol}(\Omega)$ for all sufficiently small $t$. The integral condition is its linearization by Proposition \ref{prop:cmc-first-variation}. Unless stated otherwise, we constrain the total volume of the region, not a separate volume for every boundary component. Compare \cite{ZZCMC}, Introduction, with its unrestricted $\mathcal A^c$ variations.
\end{definition}

\begin{proposition}[Realizing a zero-mean speed with an exact volume constraint]\label{prop:realize-volume-constraint}
Every smooth $f$ with $\int_\Sigma f=0$ is the initial normal speed of a smooth exactly volume-preserving variation of $\Omega$.
\end{proposition}
\begin{proof}
Use the normal tubular neighborhood of the compact embedded boundary. For small $(t,s)$, let $\Omega_{t,s}$ have boundary the normal graph of $tf+s$ over $\Sigma$. Its volume is a smooth function $V(t,s)$, with
\[
 \partial_tV(0,0)=\int_\Sigma f=0,\qquad
 \partial_sV(0,0)=\operatorname{Area}(\Sigma)>0.
\]
The implicit function theorem gives a smooth $s(t)$ with $s(0)=0$ and $V(t,s(t))=V(0,0)$. Differentiating gives $s'(0)=0$, so the initial speed is $f$. Tubular coordinates and the inverse function theorem ensure that these small normal graphs bound the corresponding smoothly varied regions. This argument uses only the first-volume-variation formula.
\end{proof}

\begin{proposition}[Fixed-volume stationarity characterizes constant mean curvature]\label{prop:fixed-volume-cmc}
A smooth compact boundary $\Sigma=\partial\Omega$ is stationary for area under all exactly volume-preserving variations if and only if its outward scalar mean curvature is one constant on all of $\Sigma$.
\end{proposition}
\begin{proof}
By the preceding proposition, stationarity gives $\int_\Sigma Hf=0$ for every smooth zero-mean function. Set $\overline H=\operatorname{Area}(\Sigma)^{-1}\int_\Sigma H$ and choose $f=H-\overline H$. Then
$0=\int H(H-\overline H)=\int(H-\overline H)^2$, so $H=\overline H$ everywhere. Conversely constant $H=c$ gives $\delta A=c\int f=0$ along every volume-preserving variation. This also shows why imposing only the total volume gives the same constant across different components: the admissible speeds allow exchange of volume between them.
\end{proof}

\begin{proposition}[Riemannian normal variation of mean curvature]\label{prop:riemannian-mean-curvature-variation}
Let a smooth two-sided hypersurface in a Riemannian manifold vary with normal velocity $f\nu$ at the time under consideration. With $H=\operatorname{div}_\Sigma\nu$,
\[
 \partial_tH=-\Delta_\Sigma f-(|A|^2+\operatorname{Ric}_M(\nu,\nu))f.
\]
Minimality is not required. We use $R(X,Y)Z=\nabla_X\nabla_YZ-\nabla_Y\nabla_XZ-\nabla_{[X,Y]}Z$, so a sectional curvature is $\langle R(e_i,\nu)\nu,e_i\rangle$. Compare \cite{HP99}, Theorem 3.2, p.~52, after replacing its inward speed by $-f$.
\end{proposition}
\begin{proof}
As in the Euclidean calculation, $\dot g_{ij}=2fh_{ij}$ and the covariant time derivative of the normal is $-\nabla f$. At a point with an orthonormal geodesic tangent frame, commute the ambient covariant derivatives in $h_{ij}=\langle\nabla_i\nu,e_j\rangle$. This gives
\[
 \dot h_{ij}=-\nabla_i\nabla_jf+fh_i{}^kh_{kj}
 -f\langle R(e_i,\nu)\nu,e_j\rangle.
\]
Indeed the commutator is $R(f\nu,e_i)\nu$, while differentiating $e_j$ contributes $f\langle\nabla_i\nu,\nabla_j\nu\rangle$. Tracing the displayed identity gives $-\Delta f+f|A|^2-f\operatorname{Ric}(\nu,\nu)$. Differentiating the inverse metric contributes $-2f|A|^2$, yielding the claim. This derives the formula independently of the second variation.
\end{proof}

\begin{proposition}[Second variation of area minus volume]\label{prop:cmc-second-variation}
If a smooth compact boundary has $H=c$, the second derivative of $\mathcal A^c=A-cV$ along a variation with initial normal speed $f$ is
\[
 Q(f,f)=\int_\Sigma\left(|\nabla f|^2-
 (|A|^2+\operatorname{Ric}_M(\nu,\nu))f^2\right)dA.
\]
Along an exactly volume-preserving variation, this is also the second derivative of area. The formula is \cite{ZZCMC}, equation (1.3), p.~6.
\end{proposition}
\begin{proof}
Differentiate $\delta\mathcal A^c=\int(H-c)f\,dA$. At the critical hypersurface, the derivatives of the speed and the area element are multiplied by $H-c=0$. The remaining term is $\int f\dot H$. Proposition \ref{prop:riemannian-mean-curvature-variation} and integration by parts give $Q(f,f)$. A tangential component of the initial velocity only differentiates the identically zero function $H-c$, so the result depends on the normal speed. If volume is constant along the variation, the second derivatives of $A$ and $A-cV$ coincide.
\end{proof}

\begin{definition}[Strong and volume-constrained CMC stability]\label{def:cmc-stability-types}
A closed CMC boundary is strongly stable if $Q(f,f)\geq0$ for every smooth $f$. It is volume-constrained stable, also called weakly stable in this setting, if the inequality holds for all $f$ with $\int_\Sigma f=0$. The latter condition is equivalent to nonnegative second variation along exactly volume-preserving paths by Proposition \ref{prop:realize-volume-constraint}. The stable $c$-hypersurfaces of \cite{ZZCMC}, Definition 2.1, use unrestricted compactly supported tests on their open domain, hence the strong notion there. Their local curvature estimates cannot be applied by silently replacing strong stability by a global volume constraint.
\end{definition}

\begin{proposition}[One scalar constraint changes index by at most one]\label{prop:cmc-constrained-index}
Let $\operatorname{ind}(Q)$ be the maximal dimension of a negative-definite subspace of smooth functions on a closed CMC boundary, and let $\operatorname{ind}_0(Q)$ use only zero-mean functions. Then
\[
 \operatorname{ind}_0(Q)\leq\operatorname{ind}(Q)
 \leq\operatorname{ind}_0(Q)+1.
\]
In particular volume-constrained stability implies unrestricted index at most one. The converse implication from index one to volume-constrained stability does not follow from this dimension estimate.
\end{proposition}
\begin{proof}
Every negative subspace in the constrained space is also negative in the full space. Conversely a finite-dimensional negative subspace $E$ intersects the kernel of the linear functional $f\mapsto\int f$ in dimension at least $\dim E-1$. The quadratic form remains negative on that intersection. Taking the supremum over $E$ proves the bounds; the argument also applies to unbounded index. This is linear algebra for the second-variation form, not an assertion about which eigenfunction has nonzero mean.
\end{proof}

\begin{proposition}[The sharp Poincare inequality on a round sphere]\label{prop:round-sphere-poincare}
On the round sphere $S_R^n$, $n\geq2$, every smooth $f$ of zero mean satisfies
\[
 \int_{S_R^n}|\nabla f|^2\,dA\geq\frac n{R^2}\int_{S_R^n}f^2\,dA.
\]
The constant is sharp, since restrictions of ambient linear coordinate functions attain equality.
\end{proposition}
\begin{proof}
The spectral theorem for the Laplace operator on a compact connected manifold gives a positive first eigenvalue $\lambda_1$ and a smooth zero-mean eigenfunction $u$, with $-\Delta u=\lambda_1u$. Integrate Bochner's identity and use $\operatorname{Ric}_{S_R^n}=(n-1)R^{-2}g$ and $|\nabla^2u|^2\geq(\Delta u)^2/n$:
\[
 0=\int\left(|\nabla^2u|^2-\lambda_1|\nabla u|^2
 +\frac{n-1}{R^2}|\nabla u|^2\right)
 \geq(n-1)\left(\frac1{R^2}-\frac{\lambda_1}{n}\right)
 \int|\nabla u|^2.
\]
Hence $\lambda_1\geq n/R^2$. Directly differentiating the coordinate functions on the sphere gives $\Delta x_i=-nx_i/R^2$, so equality is attained. The Rayleigh characterization of $\lambda_1$ gives the inequality for every zero-mean $f$. The compact spectral theorem and Bochner identity are the external analytic and differential-geometric inputs to this calculation.
\end{proof}

\begin{example}[A round sphere is weakly stable but strongly unstable]\label{ex:round-cmc-stability}
The boundary of a Euclidean ball of radius $R$ has $H=n/R$, $|A|^2=n/R^2$, and
\[
 Q(f,f)=\int_{S_R^n}\left(|\nabla f|^2-\frac n{R^2}f^2\right)dA.
\]
For $n\geq2$, Proposition \ref{prop:round-sphere-poincare} proves volume-constrained stability. The constant function gives $Q(1,1)=-nR^{-2}\operatorname{Area}(S_R^n)<0$, so the sphere is not strongly stable and its unrestricted index is exactly one by Proposition \ref{prop:cmc-constrained-index}. Coordinate-function speeds have zero quadratic form; they are the normal speeds of translations. This elementary model distinguishes the two notions of stability used in CMC problems.
\end{example}

\begin{proposition}[Strong instability under nonnegative Ricci curvature]\label{prop:cmc-ricci-strong-instability}
In an ambient manifold with $\operatorname{Ric}_M\geq0$, every nonempty closed two-sided CMC hypersurface with $c\ne0$ is unstable for unrestricted $\mathcal A^c$ variations.
\end{proposition}
\begin{proof}
The pointwise Cauchy--Schwarz inequality gives $|A|^2\geq H^2/n=c^2/n$. Therefore
$Q(1,1)=-\int(|A|^2+\operatorname{Ric}(\nu,\nu))<0$.
The constant speed changes volume, so it is not a permitted test for volume-constrained stability. This is the CMC counterpart of Corollary \ref{cor:ricci-unstable}.
\end{proof}

\begin{proposition}[Volume stability forces connectedness when the potential is positive]\label{prop:cmc-stable-connectedness}
Suppose a smooth compact CMC boundary is volume-constrained stable and
$q=|A|^2+\operatorname{Ric}_M(\nu,\nu)>0$ on every component. Then it is connected. In particular this applies to a volume-constrained stable boundary with $c\ne0$ in nonnegative Ricci curvature.
\end{proposition}
\begin{proof}
If there are two components $\Sigma_1,\Sigma_2$, choose a smooth function constant $a\ne0$ on the first, constant $b$ on the second, and zero on all others, with $a\operatorname{Area}(\Sigma_1)+b\operatorname{Area}(\Sigma_2)=0$. It has zero mean and zero gradient. Positivity of $q$ gives $Q(f,f)=-\int qf^2<0$, contradicting constrained stability. The last assertion follows from $q\geq c^2/n$. The test transfers volume between components and is admissible for a total-volume constraint.
\end{proof}

\begin{definition}[Prescribed-mean-curvature energy]\label{def:prescribed-mean-curvature-energy}
For a smooth ambient function $h$ and a smooth compact region $\Omega$, put
\[
 \mathcal A^h(\Omega)=\operatorname{Area}(\partial\Omega)-\int_\Omega h\,dV.
\]
The function $h$ is fixed in the ambient manifold as the region varies. The finite-perimeter local version was used in Proposition \ref{prop:perimeter-bulk-existence}. Constant $h=c$ gives $\mathcal A^c$; variable $h$ leads to a prescribed-mean-curvature equation, not generally to a fixed-volume constraint. Compare the constant case in \cite{ZZCMC}, equation (1.1).
\end{definition}

\begin{proposition}[First and second variation for prescribed mean curvature]\label{prop:prescribed-mean-curvature-variation}
For the outward normal convention,
\[
 \delta\mathcal A^h(f\nu)=\int_\Sigma(H-h)f\,dA.
\]
At a critical boundary $H=h|_\Sigma$, the second variation is
\[
 Q_h(f,f)=\int_\Sigma\left(|\nabla f|^2-
 (|A|^2+\operatorname{Ric}_M(\nu,\nu)+\partial_\nu h)f^2\right)dA.
\]
The normal derivative of the ambient prescribed function is part of the potential.
\end{proposition}
\begin{proof}
The transport formula for the bulk integral gives $\frac d{dt}\int_{\Omega_t}h=\int_{\Sigma_t}hf$. Subtract this from the first area variation. At $H-h=0$, differentiation leaves $\int f(\dot H-f\partial_\nu h)$. Substitute Proposition \ref{prop:riemannian-mean-curvature-variation} and integrate $-f\Delta f$ by parts. A tangential initial velocity differentiates $H-h=0$ along the hypersurface and contributes zero. The argument concerns a smooth critical boundary and supplies no regularity theorem for nonsmooth minimizers.
\end{proof}

\begin{proposition}[The CMC graph equation and oriented contact]\label{prop:cmc-graph-contact}
For a graph $z=u(x)$ with upward normal, the equation $H=c$ is
\[
 \operatorname{div}\left(\frac{Du}{\sqrt{1+|Du|^2}}\right)=-c.
\]
If two $C^2$ solutions with the same constant $c$ on a connected domain are ordered and have interior contact, they agree throughout the domain. In particular two smooth CMC hypersurfaces with the same mean-curvature constant and agreeing normals at a one-sided interior contact point agree locally. For the last assertion the ambient manifold may be Riemannian. Compare \cite{ZZCMC}, Lemma 2.7(i), p.~7.
\end{proposition}
\begin{proof}
The graph formula for $H$ follows from the normal convention in Proposition \ref{prop:graphical-mcf}. Subtracting the two equations cancels $-c$, so the difference satisfies the divergence-form linear equation in Proposition \ref{prop:graph-linearization}, with locally uniformly elliptic coefficients. Harnack applied to the nonnegative difference gives the same open-and-closed argument as for minimal graphs. In Riemannian coordinates, the difference equation has smooth bounded lower-order terms on a smaller common graph chart; the local Harnack principle gives equality near the contact point. Agreement of the normals ensures the same signed equation. Oppositely oriented touching sheets are not covered by this conclusion.
\end{proof}
